\section{Analisis horizontal}

El analisis horizontal compara una misma partida contable entre dos o mas periodos. Su finalidad es identificar aumentos, disminuciones y tendencias.

Las formulas son:

\[
\text{Variacion absoluta} = \text{Valor actual} - \text{Valor base}
\]

\[
\text{Variacion relativa} =
\frac{\text{Variacion absoluta}}{\text{Valor base}} \times 100
\]

El procedimiento consiste en seleccionar la partida, elegir un periodo base, calcular la variacion absoluta, calcular la variacion relativa e interpretar el resultado.

Este metodo responde a la pregunta: \textbf{cuanto cambio una partida entre periodos?}

\subsection{Ejemplo basico de analisis horizontal}

Enunciado: una empresa desea comparar sus ventas entre dos periodos para medir el crecimiento comercial.

Datos:

\begin{itemize}
  \item Ventas 2024: S/ 100,000.00.
  \item Ventas 2025: S/ 120,000.00.
\end{itemize}

Calculo:

\[
S/\,120,000.00 - S/\,100,000.00 = S/\,20,000.00
\]

La variacion porcentual es:

\[
\frac{S/\,20,000.00}{S/\,100,000.00} \times 100 = 20.00\%
\]

Interpretacion: las ventas aumentaron S/ 20,000.00, equivalentes al 20.00\% del periodo base.

\subsection{Ejemplo aplicado de analisis horizontal}

Enunciado: una empresa de servicios digitales contrata infraestructura cloud para alojar su plataforma. Al comparar dos periodos, la gerencia observa un incremento importante del costo mensual.

Datos:

\begin{itemize}
  \item Costo cloud 2024: S/ 8,000.00.
  \item Costo cloud 2025: S/ 12,000.00.
\end{itemize}

Calculo:

\[
S/\,12,000.00 - S/\,8,000.00 = S/\,4,000.00
\]

\[
\frac{S/\,4,000.00}{S/\,8,000.00} \times 100 = 50.00\%
\]

Interpretacion y decision: el incremento puede estar justificado si crecio la demanda de usuarios; sin embargo, si los ingresos no crecieron en la misma proporcion, debe revisarse la eficiencia de la arquitectura, proveedores o plan contratado.
